\documentclass[american]{cv}
\usepackage{pslatex}
\usepackage[T1]{fontenc}
\usepackage{geometry}
\geometry{verbose,letterpaper,tmargin=0.6in,bmargin=0.6in,lmargin=0.85in,rmargin=0.85in}
\usepackage{babel}

\makeatletter
\renewcommand{\topicmargin}{0.12\textwidth}

\begin{document}

\leftheader{Georgia Tech Research Institute\\
8800 Redstone Gateway SW Suite 100\\
Huntsville, AL 35808}

\rightheader{tel: (256) 716-2158\\
jon.fox@drfox.com\\
jon.fox@gtri.gatech.edu\\
github.com/mycr0ft}


\title{Dr. Jon R. Fox}

\maketitle
Applied Physicist with 25 years of post-doctoral experience spanning MEMS
device design and fabrication, model-based systems engineering, and signal
analysis and processing, including multiphysics modeling, hardware test, and
digital engineering environments. Programmer and data wrangler for
microcontrollers, embedded systems, automated test hardware, workstations,
and cloud-based servers.

\section{Summary of Skills}

\begin{itemize}

    \item Ph.D. in experimental condensed-matter physics.
    \item Division Chief Scientist, Architecture and System Development
        Division of the Applied Systems Laboratory (GTRI).
    \item Model-Based Systems Engineering with OMG SysML (v1 and v2),
        MagicDraw/Cameo Enterprise Architecture, and FACE conformance tooling.
    \item Open-source developer: author of SysML v2 tooling in Python
        (sysmlpy, uml2py, sysml\_canvas) with international collaborators.
    \item 25+ years of data acquisition and analysis using Python, C/C++,
        FORTRAN (77/90/95), Go, Groovy, Java, LabVIEW.
    \item Adept user (and admin) of Unix (Linux, Solaris) and Windows
        environments, office software suites, and SQL and Graph databases.

\end{itemize}

\section{Experience and Accomplishments}

\begin{topic}

\item [January 2024 - Present] Division Chief Scientist, Architecture and
System Development Division, Applied Systems Laboratory, Georgia Tech
Research Institute, Huntsville, AL.

\begin{itemize}
\item Technical direction, review, and mentoring across division programs in
model-based systems engineering, architecture, and system development.
\end{itemize}

\item [September 2019 - Present] Principal Research Scientist, Architecture
and System Development Division, Applied Systems Laboratory, Georgia Tech
Research Institute, Huntsville, AL.

\begin{itemize}
\item Model Based Systems Engineering technical analyses for US Army Aviation
Future Vertical Lift systems and software architecture.
\item GTRI Technical Lead for Future Long-Range Assault Aircraft (FLRAA)
architectural risk reduction, utilizing the Future Vertical Lift
Architectural Framework (FAF).
\item Open-source SysML v2 tooling: clean-room Python implementations of the
UML 2.5.1 metamodel (uml2py), a SysML v2 reference API and conformance
toolchain (sysmlpy, 123/123 OMG XPect conformance), an interactive
diagramming canvas (sysml\_canvas), and a SysML v2 model of the AMS
Government Reference Architecture.
\item Clean-room Python ports of the FACE Conformance Test Suite data
tooling, validated byte-for-byte against the reference Java tools.
\item Jython and Groovy scripting for data extraction and addition into
Cameo Enterprise Architecture; FACE data modeling with the FAME Awesum tool;
member of the Architecture Collaboration Working Group (ACWG).
\end{itemize}

\item [March 2010 - August 2019] System and Software Architect, Project
Manager, Applied Scientist, General Atomics Electromagnetic Systems Group,
Huntsville, AL.

\begin{itemize}
\item SysML modeling of arresting gear, electromagnetic railgun, and laser
effect systems using MagicDraw, integrated with Windchill product lifecycle
management.
\item Modeling and simulation of linear displacement transducer systems for
requirements verification and design downselect.
\item Analysis and modeling of pressure data for real-time health
prognostication.
\item Automated generation of software design documentation from SysML models
using Velocity Templating Language (VTL) and Ruby scripting.
\item Optimization of hypersonic flight trajectories with NASA's OTIS tool.
\item Information and network security analysis with CompTIA Security+
certification.
\item MEMS and nanomaterials development at the Charles Bowden Weapons
Sciences Laboratory (AMRDEC-WDI): nano-composite piezoelectric microphones,
micro-machined aluminum micro-mirrors, and MEMS health transducers; SETA
program and functional management until Mar. 2018.
\end{itemize}

\item [July 2007 - March 2010] Thin Film Process Engineer, Gulton, Inc.,
South Plainfield, NJ.

\begin{itemize}
    \item Developed thin film thermal print head (TPH) product lines from
    prototype to production; process control of existing thick film TPHs.
    \item Selection, installation, and maintenance of production physical
    vapor deposition equipment; dry and wet etch process development.
    \item Automated test equipment for resistor testing and trimming at the
    die and packaged part level; web application servers for SPC data
    collection and analysis across the factory floor.
\end{itemize}

\item [April 2000 - July 2007] Principal Scientist, Research Support
Instruments, Princeton, NJ.

\begin{itemize}
    \item Principal Investigator on SBIR/STTR grants from AFRL, ONR, NIST,
    and NASA for micro-fabricated structures: micro-machined mass
    spectrometer, energy-mass spectrograph for micro-satellites,
    micro-electroformed cross capacitor, and passive MEMS tagging.
    \item Visiting Professional Technical Staff Member II, Princeton
    University Department of Mechanical and Aerospace Engineering
    (2000--2006); Industrial Partner, Princeton Institute for the Science and
    Technology of Materials.
\end{itemize}

\item [1993 - 2000] Research Assistant, Penn State Physics Department,
University Park, PA.

\begin{itemize}
\item Analytical techniques for materials and surfaces: Raman spectroscopy,
photoluminescence, HREELS, UPS/XPS/AES, LEED/RHEED, and powder neutron
diffraction across fullerenes, Si/Ge/C nanoparticles, ferromagnetic oxides,
and II-VI epitaxial films.
\item High and ultra-high vacuum deposition: rf \& dc magnetron sputtering,
e-beam evaporation, and pulsed laser deposition.
\end{itemize}

\end{topic}

\section{Education }

\begin{topic}
\item [Ph.~D.~Physics]The Pennsylvania State University, University Park,
PA, May 2000. Thesis: ``\emph{in situ} studies of the electronic and vibrational
properties of thin films and novel materials''.
\item [B.S.~Physics]\emph{Magna~cum~laude.} University of Missouri-Rolla,
Rolla, MO, 1993. Minor in mathematics.
\end{topic}

\section{Memberships and Certifications }

\begin{itemize}
\item IEEE Computer Society; the International Council on Systems Engineering
(INCOSE).
\item INCOSE Certified Systems Engineering Professional (CSEP), March 2025.
\end{itemize}

\end{document}